\Needspace{8\baselineskip}
\Article{283}
These principles give the following method for finding all proper
representations of the binary form
\[
\varphi=ptt+2qtu+ruu
\]
of determinant $D$ by the ternary form $f$ of determinant $\Delta$.

I.\ Find all different (\textit{i.e.}, nonequivalent) values of the expression
$\surd\Delta(p,\allowbreak -q,\allowbreak r)(\mathrm{mod.}\ D)$. We solved this problem above (Article 233) for the case where $\varphi$
is a primitive form and $\Delta$ is relatively prime to $D$. The remaining
cases are very easily reduced to this one, but brevity prevents us from
explaining this in detail here. We only note that, whenever $\Delta$ is
relatively prime to $D$, the expression $\Delta(p,-q,r)$ cannot be a quadratic
residue of $D$ unless $\varphi$ is a primitive form. For if we suppose
\[
\Delta p=BB-DA',\qquad -\Delta q=BB'-DB'',\qquad \Delta r=B'B'-DA
\]
we get
\[
(DB''-\Delta q)^2=(DA'+\Delta p)(DA+\Delta r)
\]
Expanding and substituting $qq-pr$ for $D$ then gives
\[
(qq-pr)(B''B''-AA')-\Delta(Ap+2B''q+A'r)+\Delta\Delta=0
\]
From this we easily conclude that, if $p,q,r$ had a common divisor, it
would also divide $\Delta\Delta$; but then $\Delta$ could not be relatively prime
to $D$. So $p,q,r$ cannot have a common divisor; that is, $\varphi$ will be
a primitive form.

II.\ Let us denote the number of these values by $m$, and suppose
that among them there are $n$ values opposite to themselves (setting
$n=0$ when there are none). Then the remaining $m-n$ values clearly
come in opposite pairs (since we suppose that all values have been found).
Discard either value at will from each opposite pair, and in all $\frac{1}{2}(m+n)$
values will remain. Thus, \textit{for example}, of the eight values of the expression
$\surd-1(19,-3,41)$
\SourceLeaf{332}
$(\mathrm{mod.}\,770)$, namely $(39,237)$, $(171,-27)$, $(269,-83)$, $(291,-127)$, $(-39,-237)$, $(-171,27)$, $(-269,83)$,
$(-291,127)$, the last four are to be discarded as opposite to the first four.
Further, if $(B,B')$ is a value opposite to itself, clearly $2B$, $2B'$,
and therefore also $2\Delta p$, $2\Delta q$, $2\Delta r$, will be divisible by $D$.
So, if $\Delta$, $D$ are relatively prime, $2p$, $2q$, $2r$
will also be divisible by $D$. Since, by I, in this case $p$,
$q$, $r$ cannot have a common divisor either, $2$ too must
be divisible by $D$. This is impossible unless $D$ is either
$=\pm1$ or $=\pm2$. So, for all values of $D$ greater than $2$,
we will always have $n=0$ if $\Delta$ is relatively prime to $D$.

III.\ Once this is done, every proper representation of the form $\varphi$
by $f$ must clearly belong to one of the remaining values, and to
exactly one. So we must go through these values in turn and find the
representations belonging to each. To find the representations belonging
to a \textit{given} value $(B,B')$, first determine the ternary form $g=\left(\begin{smallmatrix}a,\,a',\,a''\\b,\,b',\,b''\end{smallmatrix}\right)$
whose determinant $=\Delta$ and in which $a=p$, $b''=q$, $a'=r$, $ab-b'b''=B$,
$a'b'-bb''=B'$. The values of $a''$, $b$, $b'$ are then found with the
help of the equations in II of Article 276. These readily show that, in
the case where $\Delta$, $D$ are relatively prime, $b,b'\,a''$ must be
integers (since these three numbers give integers when multiplied both
by $D$ and by $\Delta$). Now, if either one of the coefficients
$b,b',b''$ is fractional, or the forms $f$, $g$ are not equivalent,
there can be no representations of the form $\varphi$ by $f$ belonging
to $(B,B')$. But if $b,b',a''$ are integers and the forms $f$, $g$
are equivalent, every transformation of the former into the latter, such as
\[
\begin{array}{ccc}
\alpha,&\beta,&\gamma\\
\alpha',&\beta',&\gamma'\\
\alpha'',&\beta'',&\gamma''
\end{array}
\]
gives such a representation, namely
\[
x=\alpha t+\beta u,\quad x'=\alpha't+\beta'u,\quad x''=\alpha''t+\beta''u
\]
Clearly there can be no representation of this kind that cannot be
derived from some transformation. In this way, the part of the second
problem concerned with finding \textit{proper} representations has now
been reduced to the third problem.
\SourceLeaf{333}

IV.\ Further, different transformations of the form $f$ into $g$
always produce different representations, except in the single case
where the value $(B,B')$ is opposite to itself. In that case two
transformations always give a single representation. For suppose that
$f$ is also transformed into $g$ by the substitution
\[
\begin{array}{ccc}
\alpha,&\beta,&\delta\\
\alpha',&\beta',&\delta'\\
\alpha'',&\beta'',&\delta''
\end{array}
\]
(which gives the same representation as the preceding transformation).
Denoting by $k$, $\mathfrak{k}$, $\zeta$, $\eta$ the same numbers as in
II of the preceding article, we will have
\[
B=k\mathfrak{k}B+\eta\mathfrak{k}D,\quad B'=k\mathfrak{k}B'+\zeta\mathfrak{k}D
\]
So, if $k$, $\mathfrak{k}$ are supposed to be both $=+1$ or both
$=-1$, we will have (since we excluded the case $D=0$) $\zeta=0$,
$\eta=0$. This readily gives $\delta=\gamma$, $\delta'=\gamma'$, $\delta''=\gamma''$. Thus the two
transformations can differ only when one of the numbers $k$, $\mathfrak{k}$
is $+1$ and the other $-1$. Then $B\equiv-B$, $B'\equiv-B'$, $(\mathrm{mod.}\,D)$,
or the value $(B,B')$ is opposite to itself.

V.\ From the criteria for definite and indefinite forms given above
(Article 271), it readily follows that, if $\Delta$ is positive, $D$
negative, and $\varphi$ a negative form, $g$ is a negative definite
form. But if $\Delta$ is positive, and either $D$ is positive, or
$D$ is negative and $\varphi$ a positive form, $g$ is an indefinite
form. Now $f$, $g$ certainly cannot be equivalent unless they
agree in this respect. So binary forms of positive determinant, as well
as positive binary forms, clearly cannot be properly represented by a
negative ternary form; nor can negative binary forms be properly represented
by an indefinite ternary form of positive determinant. A ternary form of
the former or latter kind represents only binary forms of the latter or
former kind, respectively. In the same way, a definite ternary form of
negative determinant (\textit{i.e.}, a positive form) represents only
positive binary forms, and an indefinite one only negative binary forms
and forms of positive determinant.

\Needspace{8\baselineskip}
\Article{284}
Since the \textit{improper} representations of the binary form $\varphi$
of determinant $D$ by the ternary form $f$, whose adjoint is
$F$, are those from which improper representations
\SourceLeaf{334}
of the number $D$ by the form $F$ follow, $\varphi$ clearly cannot
be improperly represented by $f$ unless $D$ has square factors.
Suppose that all squares dividing $D$ (other than $1$) are
$ee$, $e'e'$, $e''e''$, and so on (there will be finitely many,
since we suppose that $D=0$ does not hold). Every improper representation
of the form $\varphi$ by $f$ will give a representation of the number
$D$ by $F$ in which the values of the indeterminates have one
of the numbers $e$, $e'$, $e''$, and so on as their greatest
common divisor. In this respect, for brevity, we will say that each
improper representation of the form $\varphi$ \textit{belongs} to the square
divisor $ee$ or $e'e'$ or $e''e''$, and so on. Now all representations
of the form $\varphi$ belonging to the same \textit{given} square divisor
$ee$ (whose root $e$ we take to be positive) are found by the
following rules. Their synthetic proof, which we prefer here for brevity,
will readily allow the analysis by which they were developed to be reconstructed.

\textit{First}, find all binary forms of determinant $\frac{D}{ee}$ that are
transformed into the form $\varphi$ by a proper substitution of the form
$T=\kappa t+\lambda u$, $U=\mu u$, where $T$, $U$ denote the indeterminates of
such a form; $t$, $u$ those of the form $\varphi$; $\kappa$,
$\mu$ positive integers (whose product is therefore $=e$); and
$\lambda$ a positive integer less than $\mu$ (or zero). These forms,
with their corresponding transformations, are found as follows:

Set $\kappa$ equal in turn to each positive divisor of $e$ (including
$1$ and $e$), and put $\mu=\frac{e}{\kappa}$. For each pair of fixed values
of $\kappa$, $\mu$, give $\lambda$ all integer values from $0$ to
$\mu-1$. This will certainly give all transformations. Now the form
that is transformed by any substitution $T=\kappa t+\lambda u$, $U=\mu u$ into $\varphi$
is found by determining the form into which $\varphi$ is transformed by
$t=\frac{1}{\kappa}T-\frac{\lambda}{e}U$, $u=\frac{1}{\mu}U$. In this way we obtain the forms corresponding to
each transformation; but of all these forms we must retain only those
whose three coefficients are all integers\HistoricalNote{*)}{If we could discuss this problem at greater length here, we could shorten the solution considerably. It is immediately clear that for $\kappa$ we need take only those divisors of $e$ whose square divides the first coefficient of the form $\varphi$. We reserve this problem, which also gives simpler solutions to the problems in Articles 213, 214, for another suitable occasion.}.

\textit{Second}, let $\Phi$ be one of these forms, transformed into
$\varphi$ by the substitution $T=\kappa t+\lambda u$, $U=\mu u$. Find all \textit{proper}
\SourceLeaf{335}
representations of the form $\Phi$ by $f$ (if there are any),
and express them generally by
\[
x=\mathfrak{A}T+\mathfrak{B}U,\quad x'=\mathfrak{A}'T+\mathfrak{B}'U,\quad x''=\mathfrak{A}''T+\mathfrak{B}''U\quad\ldots\quad(\mathfrak{R})
\]
Finally, from each $(\mathfrak{R})$ derive the representation
\[
x=\alpha t+\beta u,\quad x'=\alpha't+\beta'u,\quad x''=\alpha''t+\beta''u\quad\ldots\quad(\rho)
\]
by the equations
\[
\begin{array}{llll}
\alpha=\kappa\mathfrak{A},&\alpha'=\kappa\mathfrak{A}',&\alpha''=\kappa\mathfrak{A}''&\ldots\quad(R)\\
\beta=\lambda\mathfrak{A}+\mu\mathfrak{B},&\beta'=\lambda\mathfrak{A}'+\mu\mathfrak{B}',&\beta''=\lambda\mathfrak{A}''+\mu\mathfrak{B}''&
\end{array}
\]
Treat the other forms found by the first rule (if there are more) in exactly
the same way as the form $\Phi$, so that further representations are
derived from each of their proper representations. I say that this gives
all representations of the form $\varphi$ belonging to the divisor $ee$,
and each exactly once.

\textit{Proof.} I.\ That the ternary form $f$ is indeed transformed
by any substitution $(\rho)$ into $\varphi$ is so obvious that it needs
no further explanation. That each representation $(\rho)$ is improper
and belongs to the divisor $ee$ follows because the numbers $\alpha'\beta''-\alpha''\beta'$,
$\alpha''\beta-\alpha\beta''$, $\alpha\beta'-\alpha'\beta$ become $=e(\mathfrak{A}'\mathfrak{B}''-\mathfrak{A}''\mathfrak{B}')$, $e(\mathfrak{A}''\mathfrak{B}-\mathfrak{A}\mathfrak{B}'')$, $e(\mathfrak{A}\mathfrak{B}'-\mathfrak{A}'\mathfrak{B})$, respectively.
Their greatest common divisor is therefore clearly $e$ (since
$(\mathfrak{R})$ is a proper representation).

II.\ We will show that, from any given representation $(\rho)$ of the
form $\varphi$, we can find a proper representation of a form of determinant
$\frac{D}{ee}$ included among the forms found by the first rule. In other words,
the given values of $\alpha,\allowbreak \alpha',\allowbreak \alpha'',\allowbreak \beta,\allowbreak \beta',\allowbreak \beta''$ determine integer values of $\kappa$, $\lambda$,
$\mu$ satisfying the prescribed conditions, and values of $\mathfrak{A}$,
$\mathfrak{A}'$, $\mathfrak{A}''$, $\mathfrak{B}$, $\mathfrak{B}'$, $\mathfrak{B}''$ satisfying the equations
$(R)$, in exactly one way. First, the first three equations in $(R)$
immediately show that we must take $\kappa$ to be the positive greatest
common divisor of $\alpha$, $\alpha'$, $\alpha''$ (for since $\mathfrak{A}'\mathfrak{B}''-\mathfrak{A}''\mathfrak{B}'$,
$\mathfrak{A}''\mathfrak{B}-\mathfrak{A}\mathfrak{B}''$, $\mathfrak{A}\mathfrak{B}'-\mathfrak{A}'\mathfrak{B}$ must have no common divisor, $\mathfrak{A}$, $\mathfrak{A}'$,
$\mathfrak{A}''$ cannot have one either). This also determines $\mathfrak{A}$, $\mathfrak{A}'$,
$\mathfrak{A}''$, as well as $\mu=\frac{e}{\kappa}$ (which is readily seen to be an integer).
Suppose that three integers $\mathfrak{a}$, $\mathfrak{a}'$, $\mathfrak{a}''$ are chosen so
that $\mathfrak{a}\mathfrak{A}+\mathfrak{a}'\mathfrak{A}'+\mathfrak{a}''\mathfrak{A}''=1$, and for brevity write $k$ for $\mathfrak{a}\mathfrak{B}+\mathfrak{a}'\mathfrak{B}'+\mathfrak{a}''\mathfrak{B}''$.
\SourceLeaf{336}
The last three equations in $(R)$ then give $\mathfrak{a}\beta+\mathfrak{a}'\beta'+\mathfrak{a}''\beta''=\lambda+\mu k$. So there is
clearly just one value of $\lambda$ between the limits $0$ and $\mu-1$.
Since this also gives $\mathfrak{B}$, $\mathfrak{B}'$, $\mathfrak{B}''$ fixed values, all that
remains is to show that these are always integers. But we have
\[
\begin{aligned}
\mathfrak{B}=\frac{1}{\mu}(\beta-\lambda\mathfrak{A})
&=\frac{1}{\mu}(\beta(1-\mathfrak{a}\mathfrak{A})-\mathfrak{A}(\mathfrak{a}'\beta'+\mathfrak{a}''\beta''))+\mathfrak{A}k\\
&=\frac{1}{\mu}(\mathfrak{a}''(\mathfrak{A}''\beta-\mathfrak{A}\beta'')-\mathfrak{a}'(\mathfrak{A}\beta'-\mathfrak{A}'\beta))+\mathfrak{A}k\\
&=\frac{1}{e}(\mathfrak{a}''(\alpha''\beta-\alpha\beta'')-\mathfrak{a}'(\alpha\beta'-\alpha'\beta))+\mathfrak{A}k
\end{aligned}
\]
so it is clearly an integer. In the same way we readily confirm that
$\mathfrak{B}'$, $\mathfrak{B}''$ also take integer values. --- These arguments show
that there can be no improper representation of the form $\varphi$ by
$f$ belonging to the divisor $ee$ which the method given either
fails to produce or produces more than once.

If we now treat the remaining square divisors of $D$ in the same
way, finding the representations belonging to each, we will obtain all
improper representations of the form $\varphi$ by $f$.

Further, this solution readily shows that the theorem given for proper
representations at the end of the preceding article extends to improper
ones as well: in general, no positive binary form of negative determinant
can be represented by a negative ternary form, and so on. For if $\varphi$
is a binary form that cannot be properly represented by $f$ because
of that theorem, then none of the forms of determinants $\frac{D}{ee}$, $\frac{D}{e'e'}$,
and so on that contain $\varphi$ can be properly represented by $f$
either. All these forms have determinants of the same sign as $\varphi$.
Whenever these determinants are negative, they are all positive or all
negative forms according as $\varphi$ belongs to the former or the latter kind.

